% Chapter 2: Related Work

\chapter{Related Work}\label{ch:related_work}

\section{Fault-Tolerant Robot Control}\label{sec:ftc}

% Traditional FDI/FTC (Fault Detection, Isolation, and Reconfiguration)
% - Model-based approaches: analytical redundancy, observer-based
% - Data-driven approaches
% - Limitations: requires fault model, often needs downtime, not adaptive to unknown bias
% - Gap: no online learning-based approach for encoder calibration bias

\section{POMDP Solution Methods}\label{sec:pomdp_methods}

% Exact methods: belief state, PBVI, PERSEUS -> curse of dimensionality
% Approximate methods:
%   - RNN/GRU/LSTM for implicit memory -> needs multi-step history
%   - Observation augmentation -> our approach
%   - Domain randomization over hidden parameters -> related to our bias randomization
% Position our work: observation augmentation reduces POMDP to approximate MDP

\section{Robust Reinforcement Learning}\label{sec:robust_rl}

% Robust MDP / Adversarial RL
% Domain Randomization in sim-to-real transfer
%   - OpenAI: Rubik's cube with DR
%   - Tobin et al.: DR for visual transfer
% Our work: bias randomization as targeted DR on the latent variable

\section{Human-in-the-Loop Reinforcement Learning}\label{sec:hil_rl}

% Interactive Imitation Learning: DAgger
% RLPD (Ball et al.): offline demos + online interaction
% HIL-SERL (Luo et al.): human interventions for real-robot RL
% Our training framework: RLPD + keyboard intervention

\section{Safe Reinforcement Learning and Collision Avoidance}\label{sec:safe_rl}

% Constrained MDP / Safe RL
% Kiemel et al. 2024: non-negative immediate rewards for safe RL
% Backup/recovery policies in hierarchical RL
% Our backup policy: active avoidance during HIL training
